% =============================================================================
% 5.5 Trade-off Analysis
% =============================================================================

\section{Trade-off Analysis}
\label{sec:tradeoff}

The experimental results, taken together, reveal a nuanced trade-off rather than a simple win for either architecture.

\subsection{What Kubernetes Improves}
\label{subsec:k8s-improves}

\begin{description}[style=nextline]
  \item[Execution time variance.]
  Kubernetes execution time variance is lower than VM variance at all tested levels.
  The VM's variance is further inflated at L4+ by OOM-kill randomness: the Linux OOM killer selects victims non-deterministically, introducing high per-run uncertainty about whether a job will succeed at all.
  For scheduled pipelines with downstream SLA dependencies, predictable job duration is often more valuable than absolute speed.

  \item[Reliability at moderate concurrency.]
  The VM's success rate collapsed from 100\,\% to 66.7\,\% between L3 and L4, driven by memory exhaustion.
  Kubernetes maintained 100\,\% throughout, making it the only viable executor at $\geq$5 concurrent jobs for this workload and hardware configuration.

  \item[Blast radius containment at scale.]
  Kubernetes provides a \emph{hard upper bound} on memory consumption per job via cgroup-enforced pod limits.
  On the VM, OOM kills at L4+ affected multiple containers simultaneously (the OOM killer operates at the system level, not the container level), making VM blast radius uncontrolled at higher concurrency.
\end{description}

\subsection{What Kubernetes Costs}
\label{subsec:k8s-costs}

\begin{description}[style=nextline]
  \item[Container startup overhead.]
  Every Kubernetes run job incurs 19--57\,s of pre-execution overhead (pod creation + Python startup), growing from $\approx$19\,s at L1 to $\approx$57\,s at L6 under increasing concurrency.
  For the tested $\approx$63-second workload, this represents a 30--90\,\% wall-clock overhead, shrinking in relative importance for longer jobs.
  At L1--L3 where the VM is 100\% reliable, the VM is strictly faster in wall-clock time.
  From L4 onward, the VM's reliability failure makes raw speed comparison less meaningful.

  \item[Infrastructure complexity.]
  Operating a Kubernetes cluster requires significantly more expertise and operational tooling than a single VM Docker deployment.
  Helm charts, RBAC configuration, persistent volume management, and observability stack all add operational burden.
  This cost is not reflected in execution time metrics but is real and substantial for small teams.

  \item[Minimum job duration for viability.]
  To justify Kubernetes overhead, jobs should run for several minutes at minimum.
  This is a design constraint on how Dagster pipelines are structured when targeting a Kubernetes executor.
\end{description}

\subsection{Decision Framework}
\label{subsec:decision-framework}

Based on the experimental findings, a simple decision rule emerges:

\begin{itemize}
  \item \textbf{Stay on VM} if: expected peak concurrent jobs $\leq$ 3, job durations are short ($<$5 min), team prefers operational simplicity, and a 100\% success rate is achievable within the VM's RAM budget.
  \item \textbf{Migrate to Kubernetes} if: expected peak concurrent jobs $\geq$ 5, execution time variance must be controlled, job durations are $>$5 minutes (to amortise startup overhead), or per-job memory isolation is required.
  \item \textbf{The crossover threshold}: the reliability crossover (VM success rate $<$95\%) occurs at \textbf{L4 (5 concurrent jobs)} on the tested hardware. This is the empirically derived answer to the research question.
\end{itemize}

This framework answers the main research question empirically: Kubernetes becomes \emph{worth it} --- in terms of reliability, and with lower variance --- at \textbf{5 concurrent Dagster jobs} on a 4\,vCPU / 4\,GB RAM VM running a 400\,MB-per-job memory workload.

